\documentclass[]{article}
\usepackage{lmodern}
\usepackage{amssymb,amsmath}
\usepackage{ifxetex,ifluatex}
\usepackage{fixltx2e} % provides \textsubscript
\ifnum 0\ifxetex 1\fi\ifluatex 1\fi=0 % if pdftex
  \usepackage[T1]{fontenc}
  \usepackage[utf8]{inputenc}
\else % if luatex or xelatex
  \ifxetex
    \usepackage{mathspec}
  \else
    \usepackage{fontspec}
  \fi
  \defaultfontfeatures{Ligatures=TeX,Scale=MatchLowercase}
\fi
% use upquote if available, for straight quotes in verbatim environments
\IfFileExists{upquote.sty}{\usepackage{upquote}}{}
% use microtype if available
\IfFileExists{microtype.sty}{%
\usepackage{microtype}
\UseMicrotypeSet[protrusion]{basicmath} % disable protrusion for tt fonts
}{}
\usepackage[margin=1in]{geometry}
\usepackage{hyperref}
\hypersetup{unicode=true,
            pdftitle={ACTIVIDAD 1. PIRHI},
            pdfborder={0 0 0},
            breaklinks=true}
\urlstyle{same}  % don't use monospace font for urls
\usepackage{graphicx,grffile}
\makeatletter
\def\maxwidth{\ifdim\Gin@nat@width>\linewidth\linewidth\else\Gin@nat@width\fi}
\def\maxheight{\ifdim\Gin@nat@height>\textheight\textheight\else\Gin@nat@height\fi}
\makeatother
% Scale images if necessary, so that they will not overflow the page
% margins by default, and it is still possible to overwrite the defaults
% using explicit options in \includegraphics[width, height, ...]{}
\setkeys{Gin}{width=\maxwidth,height=\maxheight,keepaspectratio}
\IfFileExists{parskip.sty}{%
\usepackage{parskip}
}{% else
\setlength{\parindent}{0pt}
\setlength{\parskip}{6pt plus 2pt minus 1pt}
}
\setlength{\emergencystretch}{3em}  % prevent overfull lines
\providecommand{\tightlist}{%
  \setlength{\itemsep}{0pt}\setlength{\parskip}{0pt}}
\setcounter{secnumdepth}{5}

%%% Use protect on footnotes to avoid problems with footnotes in titles
\let\rmarkdownfootnote\footnote%
\def\footnote{\protect\rmarkdownfootnote}

%%% Change title format to be more compact
\usepackage{titling}

% Create subtitle command for use in maketitle
\newcommand{\subtitle}[1]{
  \posttitle{
    \begin{center}\large#1\end{center}
    }
}

\setlength{\droptitle}{-2em}
  \title{Diagnóstico técnico-social del protocolo de priorización de inversiones en mejoramiento de la eficiencia de conducción hídrica.}
  \pretitle{\vspace{\droptitle}\centering\huge}
  \posttitle{\par}
  \author{}
  \preauthor{}\postauthor{}
  \date{}
  \predate{}\postdate{}

\usepackage{geometry}
\geometry{letterpaper}
\usepackage{graphicx}
\usepackage{amssymb}
\usepackage{hyperref}
\usepackage{titlesec}
\usepackage{appendix}  
\usepackage{booktabs}
\usepackage{longtable}
\usepackage{multirow}
\usepackage{subcaption}
\usepackage{rotating}

\usepackage{fontspec,xltxtra,xunicode}
\defaultfontfeatures{Mapping=tex-text}
\setmainfont[		
 BoldFont={Century Gothic Bold}, 
 ItalicFont={Century Gothic Italic},
 BoldItalicFont={Century Gothic Bold Italic}
 ]{Century Gothic}

\renewcommand{\figurename}{Figura}
\renewcommand{\tablename}{Cuadro} 
\renewcommand{\contentsname}{Tabla de contenidos\\} 
\renewcommand{\listfigurename}{Índice de figuras\\}
\renewcommand{\listtablename}{Índice de cuadros\\} 
\renewcommand{\appendixname}{Anexos}
\renewcommand{\appendixtocname}{Anexos}
\renewcommand{\appendixpagename}{Anexos}


\titleformat*{\section}{\normalsize\bfseries}
\titleformat*{\subsection}{\normalsize\bfseries}
\titleformat*{\subsubsection}{\normalsize\bfseries}

\usepackage{fancyvrb}
\newcommand{\VerbBar}{|}
\newcommand{\VERB}{\Verb[commandchars=\\\{\}]}
\DefineVerbatimEnvironment{Highlighting}{Verbatim}{commandchars=\\\{\}}
% Add ',fontsize=\small' for more characters per line
\usepackage{color}
\usepackage{framed}
\usepackage{tabulary}
\definecolor{shadecolor}{RGB}{248,248,248}
\newenvironment{Shaded}{\begin{snugshade}}{\end{snugshade}}
\newcommand{\KeywordTok}[1]{\textcolor[rgb]{0.13,0.29,0.53}{\textbf{#1}}}
\newcommand{\DataTypeTok}[1]{\textcolor[rgb]{0.13,0.29,0.53}{#1}}
\newcommand{\DecValTok}[1]{\textcolor[rgb]{0.00,0.00,0.81}{#1}}
\newcommand{\BaseNTok}[1]{\textcolor[rgb]{0.00,0.00,0.81}{#1}}
\newcommand{\FloatTok}[1]{\textcolor[rgb]{0.00,0.00,0.81}{#1}}
\newcommand{\ConstantTok}[1]{\textcolor[rgb]{0.00,0.00,0.00}{#1}}
\newcommand{\CharTok}[1]{\textcolor[rgb]{0.31,0.60,0.02}{#1}}
\newcommand{\SpecialCharTok}[1]{\textcolor[rgb]{0.00,0.00,0.00}{#1}}
\newcommand{\StringTok}[1]{\textcolor[rgb]{0.31,0.60,0.02}{#1}}
\newcommand{\VerbatimStringTok}[1]{\textcolor[rgb]{0.31,0.60,0.02}{#1}}
\newcommand{\SpecialStringTok}[1]{\textcolor[rgb]{0.31,0.60,0.02}{#1}}
\newcommand{\ImportTok}[1]{#1}
\newcommand{\CommentTok}[1]{\textcolor[rgb]{0.56,0.35,0.01}{\textit{#1}}}
\newcommand{\DocumentationTok}[1]{\textcolor[rgb]{0.56,0.35,0.01}{\textbf{\textit{#1}}}}
\newcommand{\AnnotationTok}[1]{\textcolor[rgb]{0.56,0.35,0.01}{\textbf{\textit{#1}}}}
\newcommand{\CommentVarTok}[1]{\textcolor[rgb]{0.56,0.35,0.01}{\textbf{\textit{#1}}}}
\newcommand{\OtherTok}[1]{\textcolor[rgb]{0.56,0.35,0.01}{#1}}
\newcommand{\FunctionTok}[1]{\textcolor[rgb]{0.00,0.00,0.00}{#1}}
\newcommand{\VariableTok}[1]{\textcolor[rgb]{0.00,0.00,0.00}{#1}}
\newcommand{\ControlFlowTok}[1]{\textcolor[rgb]{0.13,0.29,0.53}{\textbf{#1}}}
\newcommand{\OperatorTok}[1]{\textcolor[rgb]{0.81,0.36,0.00}{\textbf{#1}}}
\newcommand{\BuiltInTok}[1]{#1}
\newcommand{\ExtensionTok}[1]{#1}
\newcommand{\PreprocessorTok}[1]{\textcolor[rgb]{0.56,0.35,0.01}{\textit{#1}}}
\newcommand{\AttributeTok}[1]{\textcolor[rgb]{0.77,0.63,0.00}{#1}}
\newcommand{\RegionMarkerTok}[1]{#1}
\newcommand{\InformationTok}[1]{\textcolor[rgb]{0.56,0.35,0.01}{\textbf{\textit{#1}}}}
\newcommand{\WarningTok}[1]{\textcolor[rgb]{0.56,0.35,0.01}{\textbf{\textit{#1}}}}
\newcommand{\AlertTok}[1]{\textcolor[rgb]{0.94,0.16,0.16}{#1}}
\newcommand{\ErrorTok}[1]{\textcolor[rgb]{0.64,0.00,0.00}{\textbf{#1}}}
\newcommand{\NormalTok}[1]{#1}
\usepackage{graphicx,grffile}
\makeatletter
\def\maxwidth{\ifdim\Gin@nat@width>\linewidth\linewidth\else\Gin@nat@width\fi}
\def\maxheight{\ifdim\Gin@nat@height>\textheight\textheight\else\Gin@nat@height\fi}
\makeatother
\usepackage{booktabs}
\usepackage{longtable}
\usepackage{array}
\usepackage{multirow}
\usepackage[table]{xcolor}
\usepackage{wrapfig}
\usepackage{float}
\usepackage{colortbl}
\usepackage{pdflscape}
\usepackage{tabu}
\usepackage{threeparttable}
\usepackage[normalem]{ulem}
\usepackage[none]{hyphenat}
\usepackage{multicol}

% Regla con color
\makeatletter
\let\old@rule\@rule
\def\@rule[#1]#2#3{\textcolor{rulecolor}{\old@rule[#1]{#2}{#3}}}
\makeatother
% Cambiar color de regla
\definecolor{rulecolor}{RGB}{183,183,183}

% Color para las letras del título
\definecolor{casinegro}{RGB}{71,71,71}

\begin{document}
\sloppy 

\newgeometry{left=3cm,bottom=0cm,top=1cm,right=3cm}

\begin{titlepage}

\begin{figure} [H]
  \includegraphics[width=.2\textwidth]{Logo/bg_gob_logo.png}
\end{figure}

\fontsize{15}{0} \selectfont{\textbf{\textcolor{casinegro}{``PLANIFICACIÓN DE LA INVERSIÓN EN REVESTIMIENTOS HÍDRICOS"}}} \vspace*{.1in}\\ % el inicio de comillas simples es ` y de doble es de ``.
\fontsize{15}{0} \selectfont{\textcolor{casinegro}{FIC-R 2016}}\\
\rule{10cm}{.05cm}
\vspace*{.1in}
\hspace*{.2in}

\fontsize{14}{0} \selectfont{\textcolor{casinegro}{Universidad de La Serena}} \vspace*{.1in}\\
\fontsize{14}{0} \selectfont{\textbf{\textcolor{casinegro}{Región de Coquimbo}}}

\vspace*{.1in}

\begin{figure}[H]
  \centering
\begin{subfigure}[c]{.2\textwidth}
  \includegraphics[width=\textwidth]{Logo/GOB_RGB.png}
\end{subfigure}
\hfill
\begin{subfigure}[c]{.23\textwidth}
  \includegraphics[width=\textwidth]{Logo/uls.png} 
\end{subfigure}
\end{figure}
\vspace*{4in}

\begin{figure}[H]
  \centering
\begin{subfigure}{.74\textwidth}
   \fontsize{14}{0} \selectfont{\textcolor{casinegro}{Marzo 2018}} \\
   {\definecolor{rulecolor}{named}{casinegro}\rule[-3mm]{10cm}{.05cm}}
\end{subfigure}
\hfill
\begin{subfigure}{.25\textwidth}
\hfill
  \includegraphics[width=\textwidth]{Logo/lgore.pdf} 
\end{subfigure}
\end{figure}


\end{titlepage}

\restoregeometry


\begin{table}[!htb]
\centering \vspace*{2.5in}
\begin{tabular}{|l|l|l|l|}
    \hline
    \multicolumn{3}{|c|}{\multirow{2}{8cm}{\centering \textbf{INFORME DE AVANCE TÉCNICO} \\
    \textbf{FIC REGIONAL} }} & \textbf{Fecha de emisión:}\\
    \cline{4-4}
    \multicolumn{3}{|c|}{\multirow{2}{8cm}} &  Abril 2018\\
    \hline
    \textbf{Fecha de inicio:}  &\textbf{Fecha de término:} &\textbf{N° de informe:} &\textbf{Duración:}\\
     20/Diciembre/2017 & 20/Octubre/2019 & 1 & 22 meses\\
    \hline
    \textbf{Nombre del proyecto:} & \multicolumn{3}{|l|}{\textbf{Planificación de la Inversión en Recursos Hídricos}}\\
    \hline
    \multirow{4}{4cm}{\textbf{Financiamiento:}} & \textbf{Monto total del proyecto:} & \multicolumn{2}{c|}{\centering M\$ 222.188}\\
    \cline{2-4}
    & \multirow{2}{4cm}{\textbf{Aporte FIC-R:}} & \multicolumn{2}{c|}{\textbf{Año 1:} M\$ 90.040}\\
    \cline{3-4}
    & & \multicolumn{2}{c|}{\textbf{Año 2:} M\$ 109.928}\\
    \cline{2-4}
    & \textbf{Aporte institución:} & \multicolumn{2}{c|}{\centering M\$ 22.220}\\
    \hline
    \textbf{Institución ejecutora:} & \multicolumn{3}{|l|}{\centering Universidad de La Serena}\\
    \hline
    \textbf{Responsable del proyecto:} & \multicolumn{3}{|l|}{\centering Dr. Pablo Álvarez Latorre}\\
    \hline
    \textbf{Cargo:} & \multicolumn{3}{|l|}{\centering Director del proyecto}\\
    \hline
    \textbf{E-mail:} & \multicolumn{3}{|l|}{\centering pabloa@userena.cl}\\
    \hline
    \textbf{Fono:} & \multicolumn{3}{|l|}{\centering 51-2554913}\\
    \hline

\end{tabular}

\end{table}

\clearpage


\clearpage
\section{Metodología}

\subsection{Componente 2. Actividad 2.1. SIG Red de distribución y conducción hídrica de Junta de Vigilancia Río Chalinga y sus Afluentes}

%incluir tabla

\begin{table}[H]
\label{my-label}
\begin{tabular}{|l|l|l|l|l|l|}
\hline
Fecha Inicio & 11-11-2017 & \multicolumn{3}{l|}{Fecha de termino}    & 1-11-2019    \\ \hline
\multicolumn{6}{|c|}{Presupuesto (M\$)}                                             \\ \hline
Total        & 1111111111 & FIC-R     & 111111111111     & ULS       & 111          \\ \hline
\multicolumn{4}{|l|}{Porcentaje de Avance}               & \multicolumn{2}{l|}{1\%} \\ \hline
\end{tabular}
\end{table}

La generación del Sistema de Información Geográfica de la red de distribución y conducción de los canales bajo el área de influencia de la Junta de Vigilancia del Río Chalinga y sus Afluentes, tiene como propósitos, la georreferenciación y caracterización de los elementos que componen a cada canal que será sometido a priorización, desde bocatoma hasta el fin del canal de distribución.

Con la generación del SIG permitirá servir de insumo para las variables que se utilizarán en el protocolo de priorización aportando información como la definición y longitud de tramos a los cuales se les evaluará las pérdidas por conducción. 

La caracterización por canal se comenzó desde bocatoma la cual tiene como principal objetivo captar y/o desviar parte o la totalidad del agua que escurre por un cauce, donde se levanto la información de los diferentes elementos que la componen, indicados a continuación:\footnote{Extraído Proyecto FIC-R Coquimbo 2013 "DISEÑO Y PUESTA EN MARCHA DE UN SISTEMA DE MONITOREO DE CAUDALES Y DE UN PROTOCOLO DE DETERMINACIÓN DE PÉRDIDAS DE AGUA, PARA LA PRIORIZACIÓN DE LAS INVERSIONES PÚBLICO-PRIVADAS" }.

\begin{itemize}
\item{\bfseries {Barrera de desviación}}: Corresponde a la estructura utilizada con el fin de desviar el agua desde la fuente (río, estero, quebrada, etc) hacia una posición en particular. Esta puede ser fabricada tanto por materiales fluviales, madera, plástico, como por hormigón o metal. Ésta barrera puede ser temporal o permanente, dependiendo del material y de la resistencia a las crecidas.
\item{\bfseries {Canal de aducción}}: Canal por donde se deriva el caudal captado.
\item{\bfseries {Compuerta de carga o punto de entrega}}: corresponde a la estructura de control en donde se entrega el recurso hídrico, este puede ser una compuerta, captación o marco partidor.
\item{\bfseries {Sección de control}}: Zona en la cual se controla el caudal entregado, en otras palabras, sección de aforo. Este puede ser de diversos tipos, dentro de los más comunes está el aforador de escurrimiento crítico y la regla limnimétrica.
\end{itemize}

Para la caracterización de los canales de distribución y conducción del recurso hídrico, se levanta información correspondiente a:
\begin{itemize}
\item{\bfseries {Revestimientos}}: Indica cuantas de las paredes del canal presenta revestimiento (1/3 una pared, 2/3 dos paredes y 3/3 100\% revestido), tipo de revestimiento (permanente o temporal), material de revestmiento, forma de la sección conductora, dimensiones de la sección conductora, todo esto indicando su longitud.
\item{\bfseries {Estructuras de distribución}}: Indica la infraestructura destinada a la entrega y gestión del recurso, ya sean compuertas, marcos partidores, captaciones con energía o gravitacionales y tacos para el primero y como gestión, compuertas frontales y laterales (compuertas de descargas).
\item{\bfseries {Observaciones}}: Indican puntos de interés, ya sean puntos críticos que puedan afectar a la conducción y distribución normal de recursos, como puntos de filtraciones, derrumbes, quebradas, embancamientos, desbordes, cruce de caminos, entre otros.
\end{itemize}


\subsubsection{Captura de datos en terreno}

Se dispuso de información geospacializada de diferentes organismos del Estado de Chile en sus diferentes formatos (shapefile, WMS, entre otros) para tener insumos base para el levantamiento de información desde terreno. 

Desde la Infraestructura de Datos Espaciales del Ministerio de Agricultura (IDE MINIAGRI), se descargó las capas vectoriales de infraestructura de riego dispuesta por la  Comisión Nacional de Riego (CNR), obteniendo la red de canales la cual considera el canal matriz, primer y segundo derivado y la capa de bocatomas la cual es informado por la Dirección General de Aguas (DGA). Al igual de la DGA se obtuvo la delimitación del área de influencia de las organizaciones de la cuenca del Choapa, desde el Sistema Integrado de Información Territorial de la Biblioteca del Congreso Nacional (SIIT-BCN) se obtuvo la red hidrográfica, con esto se pudo delimitar y tener en primera instancia los canales que serían potencialmente sometidos a priorización por organización, para este caso la cuenca del Río Chalinga.

Los datos en terreno se capturaron utilizando un GPS diferencial modelo ProXT y Juno 3B de la marca Trimble, capturando un mínimo de 30 posiciones por característica.

Cada componente de las bocatomas y  canales de distrbución, fueron caracterizados según los atributos especificados en el diccionario (VER SI GONZALO LO ANEXO PARA HACER LA REFERENCIA) de datos elaborado para la campaña de terreno, este diccionario fue ingresado en los GPS diferenciales para la captura y caracterización de los elementos.

El recorrido comenzó en las bocatomas de los canales, caracterizando las barreras de desviación, canal de aducción, compuerta de carga y sección de control. 

En el recorrido de los canales de distribución, se caracterizaron los cambios de tipos de revestimientos, las estructuras de entrega y gestión y las observaciones (puntos de interés) como: filtraciones, derrumbes, quebradas, entre otros.


\subsubsection{Posproceso de datos capturados}

Una vez obtenidos los elementos georreferenciados, se realizó la corrección diferencial de estos. Este proceso busca mejorar la calidad de los datos obtenidos en terreno, debido a que los datos brutos obtenidos del GPS varían por interferencia producto de la atmosfera y cambios en los satélites.

La corrección diferencial consiste en corregir las desviaciones del receptor móvil (GPS) utilizado, con una estación base (estática) con coordenadas conocidas. La estación base utilizada para la corrección fue SOPAC, ubicada en Santiago (33°09’01,04238’’ S, 70°40’06,80171’’ O, 723,046 msnm).

Al tomar 30 posiciones por característica georreferenciada, el promedio de estas posiciones corregidas diferencialmente incrementan, la precisión final de las características. Con la corrección diferencial se obtuvieron datos con precisiones que van entre 0,5 a 1 m en el 85\% de los casos. 





\subsubsection{Recopilación de antecedentes generales}

Se establecieron diferentes instancias para la revisión y diagnóstico del protocolo de priorización de inversiones en mejoramiento de la eficiencia de conducción hídrica. Se estableció una mesa técnica compuesta por los profesionales de las diferentes áreas del equipo, con el fin de identificar la información utilizada, metodología de recopilación y procesamiento de esta. Los resultados de esta revisión se presentan a continuación:

\paragraph{Recopilación antecedentes generales cuenca:}

Se realiza a través de una revisión bibliográfica de fuentes secundarias, como estudios de organismos referidos a los recursos hídricos de la cuenca a intervenir, como ejemplo se presentan las siguientes fuentes:

\begin{itemize}
\item Fichas de registro de bocatomas, Dirección General de Aguas
\item Canales de riego a nivel de matrices hasta 2do derivado, Comisión Nacional de Riego
\item Estudios de recursos hídricos de la zona a evaluar
\end{itemize}

\paragraph{Recopilación de antecedentes a nivel organizacional:}

Se realiza a través de reuniones y mesas de trabajo sostenidas con las organizaciones de usuarios de agua, donde se recoge información primaria la cual se detalla a continuación: 

\begin{itemize}
\item Diagrama unifilar
\item Ubicación geográfica bocatomas
\item Número de acciones totales
\item Número de acciones por comunidad de aguas
\item Dotaciones o desmarques
\item Capacidad de porteo canal
\item Número de usuarios comunidad de aguas
\item Contacto celadores
\end{itemize}

\subsubsection{Levantamiento de datos en terreno, cauce natural y artificial (canal).}

La caracterización del cauce natural, se realiza por medio de la georreferenciación y descripción de bocatomas, trazado de canales de aducción y compuertas de carga perteneciente a la comunidad de aguas. Con los puntos de trazado, se puede determinar la longitud del canal de aducción y su recorrido.

Posteriormente, en la caracterización del cauce artificial, se realiza la georreferenciación y descrpción de la infraestructura de conducción y distribución del recurso hídrico. Se capturan todas las obras de arte presentes en el recorrido del canal, desde la compuerta de carga hasta el final del recorrido, además se detalla el tipo de obra, material, finalidad, revestimeintos, además de puntos críticos como filtraciones o alteraciones visibles, quebradas, etc.

La caracterización de cauce natural y artificial se genera con el fin de proporcionar esta información a las OUA’s, como base de datos en formatos que permitan su visualización. Además con esta información se generan las variables que permitan realizar la determinación de pérdidas y posterior priorización de tramos y/o canales para el mejoramiento de la eficiencia de conducción hídrica.

\subsection{Determinación de pérdidas}

Con los insumos levantados en la caracterización, se procede a generar los tramos de aforo para la determinación de pérdidas. Estos tramos se dividen en matriz y distribución en primera instancia, para luego establecer los tramos por el cambios en el revestimiento presente, a través de aforos de los caudales de inicio y término de estos.







\subsection{Protocolo de priorización de inversiones en mejoramiento de la eficiencia de conducción hídrica}

Las variables que contiene el protocolo de priorización aplicado en SIMCA – ELQUI, las cuales se seleccionaron estadísticamente a través de la correlación de PEARSON, mayor a 0,80, en relación a la variable pérdidas generales (l/s) de los tramos.

\clearpage
\section{Resultados}

\end{document}
